\documentclass[11pt]{article}
\usepackage[T1]{fontenc}
\usepackage[utf8]{inputenc}
\usepackage{lmodern}
\usepackage{amsmath,amssymb,amsthm}
\usepackage[a4paper,margin=28mm]{geometry}
\usepackage[colorlinks=true,linkcolor=blue,citecolor=blue,urlcolor=blue]{hyperref}
\usepackage{microtype}
\newtheorem{theorem}{Theorem}[section]
\newtheorem{lemma}[theorem]{Lemma}
\newtheorem{proposition}[theorem]{Proposition}
\newtheorem{corollary}[theorem]{Corollary}
\theoremstyle{remark}
\newtheorem{remark}[theorem]{Remark}
\numberwithin{equation}{section}
\DeclareMathOperator{\co}{co_2}
\DeclareMathOperator{\ex}{ex}
\DeclareMathOperator{\exco}{exco_2}
\newcommand{\F}{K_5^{=}}
\newcommand{\Km}{K_4^{(3)-}}

% Research draft, 10 October 2026. Claims concern the specified v3 conjectures.
% No assertion of bibliographic priority or globally optimal lower bounds.
\title{Counterexamples to two squared-codegree Tur\'an conjectures}
\author{Jialin He\thanks{School of Mathematical Sciences, Key Laboratory of MEA
 (Ministry of Education), Shanghai Key Laboratory of PMMP, and Nantong Institute
 for Applied Mathematics and Artificial Intelligence, East China Normal
 University, Shanghai 200241, China. Emails: \texttt{jlhe@math.ecnu.edu.cn}.
 Supported in part by the National Natural Science Foundation of China
 (grant 12601663) and the Science and Technology Commission of Shanghai
 Municipality (grant 22DZ2229014).}}
\date{}

\begin{document}
\maketitle

\begin{abstract}
We study two five-vertex hypergraph Tur\'an problems posed by Balogh,
Clemen, and Lidick\'y. For $K_5^{=}$, the complete three-uniform hypergraph
on five vertices with two edges intersecting in one vertex deleted,
a four-part construction disproves their squared-codegree density
conjecture. It gives squared-codegree density at least
$\frac{391686088}{1029059101}\approx0.3806$ and improves the ordinary
Tur\'an lower bound to $\frac{2235+76\sqrt{114}}{5082}\approx0.5995$.
For $J_4$, the suspension of $K_4$, we correct the count for the known
iterated blow-up of the Fano complement and obtain squared-codegree
density at least $16/57$, exceeding the proposed value $7/25$.
We further prove that $16/57$ is optimal within this recursive family,
allowing arbitrary part sizes and recursion depths. The proofs use
explicit constructions, finite counting identities, and a nonnegative
weighted-square identity.
\end{abstract}

\section{Introduction}

Balogh, Clemen, and Lidick\'y~\cite{BCLpublished,BCL} study hypergraph
Tur\'an problems in which the sum of squared pair codegrees replaces the
number of edges. Their survey proposes numerical values for several
forbidden hypergraphs. We address two of its five-vertex cases: the
hypergraph $K_5^{=}$ and the suspension $J_4$ of $K_4$. For $K_5^{=}$ we
construct a counterexample to the squared-codegree conjecture and improve
the ordinary lower bound. For $J_4$ we obtain a larger lower bound by
correcting the count of an existing recursive construction.

For a positive integer $m$, write $[m]=\{1,\ldots,m\}$. For a finite
set $X$ and a nonnegative integer $r$, let $\binom Xr$ denote the family
of its $r$-element subsets.
For a finite three-uniform hypergraph $H$ with vertex set $V(H)$ and
edge set $E(H)$, write $e(H)=|E(H)|$. For distinct $u,v\in V(H)$,
let $d_H(u,v)$ be the number of edges containing $\{u,v\}$. Define
\[
 \co(H)=\sum_{\{u,v\}\in\binom{V(H)}2}d_H(u,v)^2.
\]
For a fixed three-uniform hypergraph $F$, call $H$ $F$-free if it contains
no subhypergraph isomorphic to $F$; the copy need not be induced.
Let $\ex(n,F)$ and $\exco(n,F)$ denote the
maximum values of $e(H)$ and $\co(H)$, respectively, over all $F$-free
hypergraphs on $n$ vertices. We use the normalizations
\begin{equation*}
 \pi(F)=\lim_{n\to\infty}\frac{\ex(n,F)}{\binom n3},
 \quad
 \text{and}
 \quad
 \sigma(F)=\lim_{n\to\infty}
 \frac{\exco(n,F)}{\binom n2(n-2)^2}.
\end{equation*}
Both limits exist; see~\cite[Section~1]{BCL}. For construction sequences
whose orders tend to infinity, write $n=|V(H)|$. The normalizations
$6e(H)/n^3$ and $2\co(H)/n^4$ give the ordinary and squared-codegree
limits, respectively.
We cite the published survey as~\cite{BCLpublished}
and its revised version as~\cite{BCL}. Numbered references to their
survey refer throughout to~\cite{BCL}.

\paragraph{A counterexample for $K_5^{=}$.}
Let $\F$ be the three-uniform hypergraph on $[5]$ with
\(E(\F)=\binom{[5]}3\setminus\{123,145\},
\)
where a string such as $123$ denotes $\{1,2,3\}$. Thus the two missing
edges intersect in one vertex. Balogh, Clemen, and
Lidick\'y~\cite[Conjecture~3.6]{BCL} conjecture that
\begin{equation}\label{eq:conjecture}
 \pi(\F)=\frac{227}{387},\qquad
 \sigma(\F)=\frac{2909}{8127}.
\end{equation}
Using the three-part construction described by Balogh, Clemen, and
Lidick\'y~\cite[Section~3.19]{BCL}, a machine-generated online
comment~\cite[comment~1]{ProbXiv} obtains the corrected ordinary lower bound
\[
 \pi(\F)\ge\frac59+\frac19\cdot\frac27
 =\frac{37}{63}>\frac{227}{387}.
\]
Our first result improves this corrected ordinary lower bound of $37/63$
and disproves the squared-codegree assertion in~\eqref{eq:conjecture}.

\begin{theorem}\label{thm:main}
The ordinary and squared-codegree Tur\'an densities of $\F$ satisfy
\begin{align}
 \pi(\F)&\ge\frac{2235+76\sqrt{114}}{5082}>\frac{37}{63},
           \label{eq:mainordinary}\\
 \sigma(\F)&\ge\frac{391686088}{1029059101}>\frac{2909}{8127}.
           \label{eq:mainsquared}
\end{align}
These lower bounds are approximately $0.5995$ and $0.3806$,
respectively, and are attained as limits by two sequences in the same
four-part family.
\end{theorem}

\paragraph{An improved lower bound for $J_4$.}
The hypergraph $J_4$ is obtained by adjoining one common vertex to each
edge of $K_4$; explicitly, $E(J_4)=\{123,124,125,134,135,145\}$.
Balogh, Clemen, and Lidick\'y~\cite[Conjecture~3.4]{BCL} propose
\begin{equation}\label{eq:J4conjecture}
 \sigma(J_4)=0.28=\frac7{25}.
\end{equation}
Their calculation uses balanced iterations of the complement of the
Fano plane, a construction already considered by Bollob\'as, Leader,
and Malvenuto~\cite{BLM} for the ordinary Tur\'an problem. Our second
result improves the lower bound for $\sigma(J_4)$ by correcting the
count of this recursive construction.

\begin{theorem}\label{thm:J4}
The squared-codegree Tur\'an density of $J_4$ satisfies
\begin{equation*}
 \sigma(J_4)\ge\frac{16}{57}=\frac7{25}+\frac1{1425}.
\end{equation*}
Moreover, balanced iterations of the complement of the Fano plane have limiting
squared-codegree density $16/57$.
\end{theorem}

Theorem~\ref{thm:J4} refutes~\eqref{eq:J4conjecture}. We also prove that
$16/57$ is optimal within the Fano-complement recursive family, even
with arbitrary part sizes and recursion depths
(see Lemma~\ref{lem:fano-tree}). 

Section~\ref{sec:K5} gives the four-part construction and proves
Theorem~\ref{thm:main}. In Section~\ref{sec:J4}, we prove
Theorem~\ref{thm:J4} and the subsequent optimality statement for the
Fano-complement recursive family.

\section{A counterexample for \texorpdfstring{$K_5^{=}$}{K5=}}\label{sec:K5}

We first establish a sufficient condition for a four-part construction
to be $\F$-free. We then choose the internal hypergraphs and part
proportions to obtain the two bounds in Theorem~\ref{thm:main}.

\subsection{Construction and weighted counting}\label{sec:construction}

Write $\Km$ for the complete three-uniform hypergraph on four vertices
with one edge deleted.
Let $\mathcal D$ be a loopless directed graph on $[4]$. Partition a finite
vertex set into four possibly empty parts $V_1,\ldots,V_4$.
For each $i\in[4]$, let $H_i$ be a three-uniform
hypergraph with vertex set $V_i$. Include every triple meeting three distinct parts.
For each arc $i\to j$ of $\mathcal D$, include every triple with two vertices in
$V_i$ and one in $V_j$. Finally, include the edges of each $H_i$.
Denote the resulting hypergraph by $H(\mathcal D;H_1,\ldots,H_4)$.
The use of directed graphs to prescribe edges meeting exactly two parts
also appears in Markstr\"om~\cite{Markstrom}.

The next lemma isolates the structural condition that permits internal
edges without creating the forbidden five-vertex hypergraph.

\begin{lemma}\label{lem:free}
Suppose that every $H_i$ is $\Km$-free. If every vertex of $\mathcal D$ has outdegree
at most one and $\mathcal D$ has no directed two-cycle, then
$H(\mathcal D;H_1,\ldots,H_4)$ is $\F$-free.
\end{lemma}

\begin{proof}
Suppose, for a contradiction, that the construction contains a copy of
$\F$ with vertex set $U$. Then $U$ has at most two missing triples, and
any two missing triples must intersect in exactly one vertex. Since
every four vertices of $\F$ span at least three edges and each $H_i$ is
$\Km$-free, we have $|U\cap V_i|\le3$ for every $i\in[4]$.

Suppose that $|U\cap V_i|=3$ for some $i$. For any
$v\in U\setminus V_i$, say $v\in V_j$, the arc $i\to j$ must be present. Otherwise the three triples consisting of $v$ and two vertices of
$U\cap V_i$ are missing. As $i$ has outdegree at most one, both vertices
of $U\setminus V_i$ lie in the same part $V_j$. The absence of directed
two-cycles implies that $j\to i$ is absent. These two vertices, together
with each vertex of $U\cap V_i$, give three missing triples, a
contradiction. Thus $|U\cap V_i|\le2$ for every $i\in[4]$.

Since $|U|=5$ and there are only four parts, some part satisfies
$U\cap V_i=\{a,b\}$. Among the three triples $\{a,b,x\}$ with
$x\in U\setminus V_i$, at most one can be missing, since any two
intersect in the pair $\{a,b\}$. Thus there are distinct
$c,d\in U\setminus V_i$ such that both $\{a,b,c\}$ and $\{a,b,d\}$
are edges. By the construction, the parts containing $c$ and $d$ must
have labels that are outneighbors of $i$. Since $i$ has outdegree at
most one, $c$ and $d$ lie in the same part $V_j$, with $i\to j$.
The absence of directed two-cycles implies that $j\to i$ is absent.
Hence $\{a,c,d\}$ and $\{b,c,d\}$ are missing and intersect in the
pair $\{c,d\}$, a contradiction. Therefore the construction is
$\F$-free.
\end{proof}

We now specify the internal hypergraphs using a classical six-vertex
construction. To make the counting self-contained, define $S$ on $[6]$ by
\begin{equation*}
 E(S)=\{123,125,136,145,146,234,246,256,345,356\}.
\end{equation*}
This is the hypergraph used in the iterated construction of
Frankl and F\"uredi~\cite{FF}; see also~\cite[Section~3.9]{BCL}.
The edge list shows that every four vertices of $S$ span exactly two
edges and every pair has codegree two.

We define a sequence $(S_k)_{k\ge0}$ recursively. Let $S_0$ consist of
one vertex and no edges. For each $k\ge1$, having defined $S_{k-1}$,
take six vertex-disjoint copies of $S_{k-1}$ with vertex sets
$W_1,\ldots,W_6$. To form $S_k$, retain all edges within these copies
and, for each $\{i,j,\ell\}\in E(S)$, add every triple with one vertex
in each of $W_i,W_j,W_\ell$. No other edges are added. In particular,
no edge meets exactly two of the sets $W_i$. Thus $S_1\cong S$ and
$|V(S_k)|=6|V(S_{k-1})|=6^k$ for $k\ge1$.

\begin{lemma}\label{lem:internal}
For every integer $k\ge0$, the hypergraph $S_k$ is $\Km$-free and, with
$n=6^k$ vertices, satisfies
\begin{equation}\label{eq:internalfinite}
 e(S_k)=\frac{n^3-n}{21},\quad
 \text{and} 
 \quad
 \co(S_k)=\frac{2(n^4-n)}{43}.
\end{equation}
As $k\to\infty$, the ordinary and squared-codegree densities of $S_k$
tend to $\rho=2/7$ and $\tau=4/43$, respectively.
\end{lemma}

\begin{proof}
We proceed by induction on $k$. The case $k=0$ is immediate. Suppose
that $k\ge1$ and the assertions hold for $S_{k-1}$.
Let $U$ be a four-element subset of $V(S_k)$. If $U\subseteq W_i$ for
some $i$, it spans at most two edges by induction. If the vertices of
$U$ lie in four distinct parts, they span exactly two edges by the
four-vertex property of $S$. Otherwise, some part satisfies
$r=|U\cap W_i|\in\{2,3\}$. Every triple consisting of two vertices
of $U\cap W_i$ and one vertex of $U\setminus W_i$ is missing. There
are $\binom r2(4-r)\ge2$ such triples, so $U$ spans at most two edges.
Hence $S_k$ is $\Km$-free.

For the counts, put $m=6^{k-1}$, so $|W_i|=m$ and $n=6m$.
The ten edges of $S$ contribute $10m^3$ edges meeting three parts.
Since every pair of vertices of $S$ has codegree two, each pair in
distinct parts of $S_k$ has codegree $2m$. There are $15m^2$ such
pairs, contributing $15m^2(2m)^2=60m^4$ to $\co(S_k)$.
Pairs within a part have only internal common neighbors. Consequently,
\[
 e(S_k)=10m^3+6e(S_{k-1}),\quad
 \text{and}
 \quad
 \co(S_k)=60m^4+6\co(S_{k-1}).
\]
Applying the induction hypotheses to $S_{k-1}$, which has $m$ vertices,
and using $n=6m$, we obtain
\begin{align*}
 e(S_k)&=10m^3+\frac{6(m^3-m)}{21}
        =\frac{(6m)^3-6m}{21}
        =\frac{n^3-n}{21},\\
 \intertext{and}
 \co(S_k)&=60m^4+\frac{12(m^4-m)}{43}
         =\frac{2((6m)^4-6m)}{43}
         =\frac{2(n^4-n)}{43}.
\end{align*}
This proves~\eqref{eq:internalfinite} and completes the induction.
Multiplying the first by $6/n^3$ and the second by $2/n^4$ gives
the limits $2/7$ and $4/43$.
\end{proof}

We now count hypergraphs of the form $H(\mathcal D;H_1,\ldots,H_4)$.
Fix a loopless directed graph $\mathcal D$ on $[4]$ with outdegree at
most one and no directed two-cycle. Let $x=(x_1,\ldots,x_4)$ be a vector
of nonnegative target part proportions with $\sum_i x_i=1$.
All class indices in the following definitions and weighted formulas
range over $[4]$. Define
\[
 L_i=\sum_{j:i\to j}x_j,\qquad
 C_{ij}=1-x_i-x_j+\mathbf1_{i\to j}x_i+\mathbf1_{j\to i}x_j
 \quad(i\ne j).
\]
Here $\mathbf1_{i\to j}$ is $1$ if $i\to j$ is an arc of $\mathcal D$
and $0$ otherwise. Sums indexed by $i\to j$ run over the arcs of
$\mathcal D$.
For sequences in this family with orders tending to infinity and
part proportions tending to $x$, $L_i$ is the limiting proportion of
common neighbors outside $V_i$ for a pair
within $V_i$, while $C_{ij}$ is the limiting proportion of all common
neighbors for a pair in distinct parts $V_i,V_j$. Both proportions are
taken relative to the total number of vertices.

\begin{lemma}\label{lem:weighted}
There is a sequence of $\F$-free hypergraphs of the form
$H(\mathcal D;H_1,\ldots,H_4)$ whose orders tend to infinity and whose
part proportions tend to $x$. Their ordinary and squared-codegree
densities converge, respectively, to
\begin{align}
 P_{\mathcal D}(x)&=6\sum_{i<j<\ell}x_ix_jx_\ell
          +3\sum_{i\to j}x_i^2x_j+\rho\sum_i x_i^3,
          \label{eq:weightedP}\\
 Q_{\mathcal D}(x)&=2\sum_{i<j}x_ix_j C_{ij}^2
          +\sum_i x_i^2\bigl(L_i^2+2\rho x_iL_i+\tau x_i^2\bigr),
          \label{eq:weightedQ}
\end{align}
where $\rho=2/7$ and $\tau=4/43$.
\end{lemma}

\begin{proof}
We form $H(\mathcal D;H_1,\ldots,H_4)$ by choosing each internal
hypergraph $H_i$ as a suitable blow-up of $S_k$ from
Lemma~\ref{lem:internal}.

First suppose that $x_1,\ldots,x_4$ are rational. For each integer
$k\ge0$, let $n$ be $6^k$ times a fixed positive common denominator of
these four numbers. Thus $n\to\infty$ as $k\to\infty$.
For each $i\in[4]$, replace each vertex of a separate copy of $S_k$
by $nx_i/6^k$ independent clones, including all triples whose original
vertices form an edge. Denote the resulting internal hypergraph by
$H_i$ and its vertex set by $V_i$. Since $S_k$ has $6^k$ vertices,
$|V_i|=nx_i$, with $V_i$ empty when $x_i=0$. 

Four clones of distinct original vertices span at most two edges by
Lemma~\ref{lem:internal}. If two of the four clones have the same original
vertex, both triples containing that pair are absent, so again at most
two edges remain. Thus each $H_i$ is $\Km$-free, and Lemma~\ref{lem:free}
shows that $H=H(\mathcal D;H_1,\ldots,H_4)$ is $\F$-free.

Cloning multiplies the edge count by the cube of the number of clones
per vertex, and the squared-codegree sum by its fourth power since a pair
of distinct original vertices gives the square of that number of clone
pairs, each with internal codegree multiplied by the same number.
Two clones of the same vertex have codegree zero in $H_i$.
Consequently,~\eqref{eq:internalfinite} gives
\begin{align*}
 \frac{6e(H_i)}{n^3}
   &=\frac{6x_i^3}{6^{3k}}e(S_k)
     =\rho x_i^3(1-6^{-2k}),\quad
    \text{and}\quad
 \frac{2\co(H_i)}{n^4}
   =\frac{2x_i^4}{6^{4k}}\co(S_k)
     =\tau x_i^4(1-6^{-3k}).
\end{align*}
Counting edges that meet three parts, two parts, or one part, and using
$|V_i|=nx_i$, we obtain the exact identity
\[
\begin{aligned}
 \frac{6e(H)}{n^3}
   ={}&6\sum_{i<j<\ell}x_ix_jx_\ell
       +3\sum_{i\to j}x_i\left(x_i-\frac1n\right)x_j
    +\sum_i\frac{6e(H_i)}{n^3}.
\end{aligned}
\]
Substituting the first internal count above and letting $k\to\infty$
gives~\eqref{eq:weightedP}.

To count squared codegrees, let
$\varepsilon_{ij}=\mathbf1_{i\to j}+\mathbf1_{j\to i}$.
For $u\in V_i$ and $v\in V_j$ with $i\ne j$, the codegree is
\[
\begin{aligned}
 d_H(u,v)&=n(1-x_i-x_j)
       +\mathbf1_{i\to j}(nx_i-1)+\mathbf1_{j\to i}(nx_j-1)
        =nC_{ij}-\varepsilon_{ij}.
\end{aligned}
\]
The terms $nx_i-1$ and $nx_j-1$ exclude the endpoints $u$ and $v$,
respectively. For distinct $u,v\in V_i$, we have
$d_H(u,v)=nL_i+d_{H_i}(u,v)$.
Expanding the square gives
\[
 d_H(u,v)^2=n^2L_i^2+2nL_i\,d_{H_i}(u,v)+d_{H_i}(u,v)^2.
\]
When we sum over all pairs in $V_i$, the internal codegrees sum to
$3e(H_i)$, because each internal edge contains three pairs. Hence the
middle terms contribute $2nL_i\cdot3e(H_i)=6nL_i e(H_i)$, while the
last terms contribute $\co(H_i)$.
There are $n^2x_ix_j$ pairs between $V_i$
and $V_j$, and $\binom{nx_i}{2}$ pairs within $V_i$. Summing their
squared codegrees and multiplying by $2/n^4$ gives
\begin{equation*}
\begin{split}
 \frac{2\co(H)}{n^4}
   ={}&2\sum_{i<j}x_ix_j\left(C_{ij}-\frac{\varepsilon_{ij}}n\right)^2
     +\sum_i\left[
       x_i\left(x_i-\frac1n\right)L_i^2
       +2L_i\frac{6e(H_i)}{n^3}
       +\frac{2\co(H_i)}{n^4}\right].
\end{split}
\end{equation*}
Substituting the two internal counts above and letting $k\to\infty$,
the three terms in brackets tend to $x_i^2L_i^2$,
$2\rho x_i^3L_i$, and $\tau x_i^4$, respectively. The first sum tends
to $2\sum_{i<j}x_ix_jC_{ij}^2$, proving~\eqref{eq:weightedQ}.

Finally, we extend the construction to real part proportions; this
allows us to use the irrational proportions in the proof of
Theorem~\ref{thm:main}. Given any such $x$, choose a sequence
$(x^{(j)})_{j\ge1}$ of rational four-dimensional vectors with nonnegative
coordinates summing to one and $x^{(j)}\to x$.
For each integer $j\ge1$, the rational case provides an $\F$-free
hypergraph $H$ on $n\ge j$ vertices with part proportions $x^{(j)}$,
such that its two normalized densities differ from
$P_{\mathcal D}(x^{(j)})$ and $Q_{\mathcal D}(x^{(j)})$, respectively,
by less than $1/j$. For these chosen graphs,
\(
 \left|\frac{6e(H)}{n^3}-P_{\mathcal D}(x)\right|
\le\frac1j+\left|P_{\mathcal D}(x^{(j)})-P_{\mathcal D}(x)\right|,
 \)
 and
 \(
 \left|\frac{2\co(H)}{n^4}-Q_{\mathcal D}(x)\right|
\le\frac1j+\left|Q_{\mathcal D}(x^{(j)})-Q_{\mathcal D}(x)\right|.
\)
Both right-hand sides tend to zero by continuity of the two polynomials.
Since $n\ge j$ and $x^{(j)}\to x$, these graphs also have orders tending
to infinity and part proportions tending to $x$, as required.
\end{proof}

\subsection{The two density bounds}

\begin{proof}[Proof of Theorem~\ref{thm:main}]
We use the same four-part family for both bounds, with different part
proportions. Relabel the four parts $1,2,3,4$ as $A,B,C,D$,
respectively, and write $x=(x_A,x_B,x_C,x_D)$. Let $\mathcal D$ have
the arcs
\begin{equation*}
 A\to B,\quad B\to C,\quad C\to A,\quad D\to A.
\end{equation*}
Every vertex of $\mathcal D$ has outdegree one, and there is no directed
two-cycle. Thus Lemma~\ref{lem:weighted} provides an $\F$-free sequence
with limiting densities $P_{\mathcal D}(x)$ and $Q_{\mathcal D}(x)$ for
every vector $x$ of nonnegative part proportions summing to one.
In particular, $\pi(\F)\ge P_{\mathcal D}(x)$ and
$\sigma(\F)\ge Q_{\mathcal D}(x)$. It remains to choose suitable
proportions for each bound.

For the ordinary density, let $t\in[0,1]$ be a parameter to be chosen
later, and set $x_A=x_B=x_C=(1-t)/3$ and $x_D=t$.
Writing $p(t)=P_{\mathcal D}(x)$, substitution
in~\eqref{eq:weightedP} gives
\begin{equation*}
\begin{split}
 p(t)=\frac{37}{63}(1-t)^3+2t(1-t)^2+t^2(1-t)+\frac27t^3
     =\frac{37+15t-78t^2+44t^3}{63}.
\end{split}
\end{equation*}
Its derivative is $(15-156t+132t^2)/63$. The two roots are
$(13\pm\sqrt{114})/22$, and only the smaller root lies in $[0,1]$.
On $[0,1]$, the derivative is positive before that root and negative
after it.
Therefore
\begin{equation}\label{eq:pmax}
 \max_{0\le t\le1}p(t)
 =p\left(\frac{13-\sqrt{114}}{22}\right)
 =\frac{2235+76\sqrt{114}}{5082}.
\end{equation}
This proves~\eqref{eq:mainordinary}. Since $p(0)=37/63$ and $p'(0)>0$,
the bound strictly exceeds the corrected three-part value. The maximum
in~\eqref{eq:pmax} concerns only the indicated one-parameter subfamily.

For the squared-codegree density, take
$x=\left(\frac{13}{43},\frac{12}{43},\frac{12}{43},\frac{6}{43}\right)$ in the
order $A,B,C,D$. Then
\(
 (L_A,L_B,L_C,L_D)
 =\left(\frac{12}{43},\frac{12}{43},\frac{13}{43},\frac{13}{43}\right),\)
 and
\[ (C_{AB},C_{AC},C_{AD},C_{BC},C_{BD},C_{CD})
 =\left(\frac{31}{43},\frac{30}{43},\frac{30}{43},
          \frac{31}{43},\frac{25}{43},\frac{25}{43}\right).
\]
Substitution in~\eqref{eq:weightedQ} gives
\begin{equation*}
 \begin{split}
 Q_{\mathcal D}(x)
 =\frac{1}{43^4}\left(1253292+\frac{4\cdot72372}{7}
                      +\frac{4\cdot71329}{43}\right)
 =\frac{391686088}{1029059101}
  =\frac{2909}{8127}+\frac{90033181}{3969227961}.
 \end{split}
\end{equation*}
Together with Lemma~\ref{lem:weighted}, this proves
\eqref{eq:mainsquared} and completes the proof.
\end{proof}

The two lower bounds use different proportions. Neither the full
four-variable optimization nor the exact values of $\pi(\F)$ and
$\sigma(\F)$ are determined here.

\section{An improved lower bound for \texorpdfstring{$J_4$}{J4}}\label{sec:J4}

The link of a vertex $v$ in a three-uniform hypergraph is the graph on
the remaining vertices in which $uw$ is an edge exactly when $vuw$ is
a hyperedge. A hypergraph is $J_4$-free precisely when all its vertex
links are $K_4$-free. We use this characterization to verify the Fano
construction and then count its squared codegrees.

\begin{proof}[Proof of Theorem~\ref{thm:J4}]
Let $B$ be the complement of a Fano plane on vertex set $[7]$.
The triples of the plane are called its lines; these are precisely
the triples missing from $B$. The hypergraph $B$ has
$28$ edges; each pair has codegree four. At a fixed vertex, its link
is $K_6$ minus a perfect matching, that is, $K_{2,2,2}$.
Set $B_0$ to be a single vertex with no edges. For each integer $k\ge1$,
construct $B_k$ from seven vertex-disjoint copies of $B_{k-1}$ labelled
by the vertices of $B$. Retain the edges within each copy, and add every
triple meeting three distinct copies whose labels form an edge of $B$.
No triples meeting exactly two copies are added.
In the link of a vertex of $B_k$, the portion outside its
first-level part is a blow-up of $K_{2,2,2}$, hence a complete
tripartite graph. The portion inside
that part is a vertex link in $B_{k-1}$, and there are no link edges
between these two portions. Induction shows that the entire link has
no $K_4$. Thus $B_k$ is $J_4$-free.

For $k\ge1$, put $n=7^k$ and $m=7^{k-1}$. Pairs in distinct first-level parts
have codegree $4m$. There are $21m^2$ such pairs, contributing
$21m^2(4m)^2=336m^4$ to $\co(B_k)$.
Pairs within one part have no common neighbors outside that part,
since no edge meets exactly two parts. Thus each of the seven parts
contributes precisely $\co(B_{k-1})$, and
\[
 \co(B_k)=336\cdot7^{4(k-1)}+7\co(B_{k-1}).
\]
Dividing the recurrence by $7^{4k}$ gives a factor $7^{-3}$ in front
of the preceding normalized term. Starting from $\co(B_0)=0$, we obtain
\[
 \frac{\co(B_k)}{7^{4k}}
 =\frac{336}{7^4}\sum_{j=0}^{k-1}7^{-3j}
 =\frac8{57}(1-7^{-3k}),
\]
where the last equality uses the geometric sum and
$336/[7^4(1-7^{-3})]=8/57$. Since $n=7^k$, this gives the exact count
\begin{equation}\label{eq:J4finite}
 |V(B_k)|=n=7^k,\quad
 \text{and}\quad \co(B_k)=\frac{8(n^4-n)}{57}.
\end{equation}
Consequently,
\[
 \frac{2\co(B_k)}{n^4}=\frac{16}{57}(1-n^{-3})
 \longrightarrow\frac{16}{57}.
\]
Since every $B_k$ is $J_4$-free, this limit proves the asserted lower bound.
\end{proof}

The factor seven in the recurrence accounts for all seven recursive
parts. After normalization by $n^4$, successive levels therefore
contribute in the ratio $7/7^4=1/343$.
The calculation in~\cite[Section~3.15]{BCL} uses $1/7^4$ instead.
Correcting this factor gives $16/57$ and refutes the numerical equality
$\sigma(J_4)=0.28$ in Conjecture~3.4 of that version, using the same
underlying construction.

\paragraph{Optimality within the recursive family.}
We next show that $16/57$ remains optimal when class sizes and
recursion depths vary. Let $B$ again denote the Fano complement.
A \emph{Fano-complement recursive construction} is specified by a
finite rooted tree with independent vertex classes at its leaves.
Each internal node partitions its vertex set into at most seven
children labelled by distinct vertices of $B$. Include all triples
meeting three children whose labels form an edge of $B$, and no triples
meeting exactly two children. Edges within each child are defined
recursively. This is a special case of the recursive patterns
of Liu and Pikhurko~\cite[Definition~1.1]{LP}.

For a real weight vector $x=(x_1,\ldots,x_7)$ on the vertices of $B$, put
$W=\sum_{i=1}^7 x_i$, $d_{ij}=\sum_{k:\,ijk\in E(B)}x_k$ for distinct
$i,j\in[7]$, and
$N(x)=2\sum_{i<j}x_ix_jd_{ij}^2$.
The weights need not sum to one.
When the weights are class sizes, $d_{ij}$ is the codegree of a pair
between classes $i,j$, and $N(x)/2$ is the total squared-codegree
contribution from pairs in distinct classes.
The following lemma bounds this contribution.

\begin{lemma}\label{lem:fano-weight-sos}
For nonnegative weights,
\(
 N(x)\le\frac{16}{57}\left(W^4-\sum_i x_i^4\right).
\)
\end{lemma}

\begin{proof}
For each $1\le i<j\le7$, let $k$ be the third vertex on their Fano line.
Label the remaining vertices $a,b,c,d$ so that $kab$ and $kcd$ are lines.
Then $d_{ij}=x_a+x_b+x_c+x_d$. For all real weights,
$16\left(W^4-\sum_i x_i^4\right)-57N(x)$ equals
\begin{equation*}
\sum_{i<j}(x_i-x_j)^2\Bigl[
      12d_{ij}^2+2(8x_i+8x_j+x_k)d_{ij}+24(x_a+x_b)(x_c+x_d)+58(x_ax_b+x_cx_d)\Bigr].
\end{equation*}
To check the identity, expand the bracket for each pair $i,j$.
For $u\in\{a,b,c,d\}$, the coefficients of
$x_u^2,x_ix_u,x_jx_u,x_kx_u$ are $12,16,16,2$, respectively.
The coefficients of $x_ax_b$ and $x_cx_d$ are $24+58=82$;
the other mixed products among $x_a,x_b,x_c,x_d$ have coefficient
$24+24=48$. Multiplication by $x_i^2+x_j^2-2x_ix_j$ and summation
over pairs give the right-hand coefficients in the table below,
where $p,q,r,s\in[7]$ are distinct.
\begin{center}
\begin{tabular}{llll}
\hline
Monomial & Fano condition & Left-hand side & Right-hand side \\
\hline
$x_p^4$ & --- & $0$ & $0$ \\
$x_p^3x_q$ & --- & $64$ & $4\cdot16=64$ \\
$x_p^2x_q^2$ & --- & $96$ & $2\cdot4\cdot12=96$ \\
$x_p^2x_qx_r$ & $pqr$ is a line & $192$ & $4\cdot48=192$ \\
$x_p^2x_qx_r$ & $pqr$ is not a line & $78$ & $78$ \\
$x_px_qx_rx_s$ & contains a line & $-300$ & $-300$ \\
$x_px_qx_rx_s$ & contains no line & $-984$ & $-984$ \\
\hline
\end{tabular}
\end{center}
If $pqr$ is not a Fano line, the $x_p^2$ terms from the six summands
indexed by pairs containing $p$ contribute
$2\cdot16+2\cdot2+48+82$ to the coefficient of $x_p^2x_qx_r$.
The mixed terms from the pairs $\{p,q\}$, $\{p,r\}$, and $\{q,r\}$
contribute $-2(16+16+12)$. Their sum is $78$, matching
$16\cdot12-57\cdot2$ on the left.
For a monomial with four distinct indices, only mixed terms contribute
on the right: its coefficient is $-2(3\cdot2+3\cdot48)=-300$ if
the four points contain a Fano line, and $-2\cdot6\cdot82=-984$
otherwise. On the left, the corresponding coefficients of $N(x)$
are $4\cdot3=12$ and $4\cdot6=24$, giving
$16\cdot24-57\cdot12=-300$ and $16\cdot24-57\cdot24=-984$.
The table covers all monomials of degree four, proving the identity.
For nonnegative weights, its right-hand side is nonnegative, and the
inequality follows.
\end{proof}

\begin{lemma}\label{lem:fano-tree}
Every Fano-complement recursive construction $H$ on $n\ge1$ vertices is
$J_4$-free and satisfies
\[
 \co(H)\le\frac8{57}(n^4-n).
\]
The balanced iterations on $7^k$ vertices attain equality for every
integer $k\ge0$. Consequently, the supremum of the limiting
squared-codegree densities in this family is $16/57$, even when class
sizes and recursion depths are chosen independently at different nodes.
\end{lemma}

\begin{proof}
We induct on the height of the recursion tree.
If the root is a leaf, then $H$ is independent, so it is
$J_4$-free and $\co(H)=0\le8(n^4-n)/57$.
Otherwise, let $H_1,\ldots,H_7$ be the hypergraphs in the root's
children, using a hypergraph with no vertices for each unused label. Put
$n_i=|V(H_i)|$, so that $n=\sum_i n_i$.
As in the proof of Theorem~\ref{thm:J4}, each vertex link is the
disjoint union of its link within its child class and a complete
tripartite graph, possibly with empty parts. Both are $K_4$-free by
induction and construction, respectively, so $H$ is $J_4$-free.

There are $n_in_j$ pairs between children $i,j$,
each with codegree $\sum_{k:\,ijk\in E(B)}n_k$.
Pairs within one child have no common neighbors outside that child,
since no edge meets exactly two children. Hence
\[
 \co(H)=\sum_i\co(H_i)
   +\sum_{i<j}n_in_j\left(\sum_{k:\,ijk\in E(B)}n_k\right)^2
 =\sum_i\co(H_i)+\frac12N(n_1,\ldots,n_7).
\]
Here $N(n_1,\ldots,n_7)$ is $N(x)$ evaluated at $x_i=n_i$;
its half is exactly the contribution from pairs in distinct children.
Since $W=\sum_i n_i=n$, Lemma~\ref{lem:fano-weight-sos} bounds this
term. The induction hypothesis bounds
$\co(H_i)$ for each nonempty child; empty children contribute zero.
Thus
\[
 \co(H)\le\frac8{57}\left(n^4-\sum_i n_i^4\right)
            +\frac8{57}\sum_i(n_i^4-n_i)
          =\frac8{57}(n^4-n).
\]
The last equality uses the cancellation of the fourth-power terms
and $\sum_i n_i=n$. Hence
$2\co(H)/n^4\le(16/57)(1-n^{-3})$, giving limiting
squared-codegree density at most $16/57$.
By~\eqref{eq:J4finite}, the balanced iterations attain equality in
the finite bound and have densities tending to $16/57$.
\end{proof}

Theorem~\ref{thm:J4} and Lemma~\ref{lem:fano-tree} support
$\sigma(J_4)=16/57$ as a natural numerical revision of the conjecture
in~\cite[Conjecture~3.4]{BCL}. This retains the proposed extremal
construction with its corrected density. 

\section*{Declaration of AI use}

The author used GPT-6 to assist with exploring proof strategies,
checking proofs, and improving the exposition. All mathematical
arguments, results, and conclusions were reviewed and verified by the
author, who takes full responsibility for the content of this paper.

\begin{thebibliography}{9}

\bibitem{BCLpublished}
J.~Balogh, F.~C.~Clemen, and B.~Lidick\'y,
Hypergraph Tur\'an problems in $\ell_2$-norm,
in \emph{Surveys in Combinatorics 2022},
London Math. Soc. Lecture Note Ser., vol.~481,
Cambridge University Press, Cambridge, 2022, pp.~21--63.

\bibitem{BCL}
J.~Balogh, F.~C.~Clemen, and B.~Lidick\'y,
Hypergraph Tur\'an problems in $\ell_2$-norm,
arXiv:2108.10406v3,
2025.

\bibitem{BLM}
B.~Bollob\'as, I.~Leader, and C.~Malvenuto,
Daisies and other Tur\'an problems,
\emph{Combin. Probab. Comput.} \textbf{20} (2011), 743--747.

\bibitem{FF}
P.~Frankl and Z.~F\"uredi,
An exact result for 3-graphs,
\emph{Discrete Math.} \textbf{50} (1984), 323--328.

\bibitem{LP}
X.~Liu and O.~Pikhurko,
Finite hypergraph families with rich extremal Tur\'an constructions via mixing patterns,
\emph{Forum Math. Sigma} \textbf{13} (2025), e53.

\bibitem{Markstrom}
K.~Markstr\"om,
Two questions of Erd\H{o}s on hypergraphs above the Tur\'an threshold,
\emph{J. Graph Theory} \textbf{76} (2014), 101--105.

\bibitem{ProbXiv}
ProbXiv,
Comment~1 on ``Hypergraph Tur\'an problems in $\ell_2$-norm'',
August~17, 2026. Machine-generated online contribution
(labelled GPT-5.5 xhigh); the page reports no human review.
Accessed October~11, 2026.
\url{https://probxiv.com/problems/hypergraph-turan-problems-in-ell-2-norm-7}.

\end{thebibliography}
\end{document}
